\section{The consistency-tolerance theorem}
\label{sec:theorem}

This section states and proves the first contribution: a gauge lemma, a measured
affine reducible/irreducible split with a proven removability structure, a closed-form
consistency tolerance, and the design-law corollary. Each result carries an explicit
honesty label. The algebraic statements (Lemma~\ref{lem:gauge}, the removability
structure of Lemma~\ref{lem:split}, the triangle bound, and the $\tau(B)$ inversion
form) are \Proven{}; the affine \emph{attribution} of the recovered rotation into a
frame-tie and a Moon-excess is a modelling interpretation, stated as such
(Remark~\ref{rem:attribution}); and the $B_{\mathrm{eff}}$ inflation and every
system-level magnitude are \Modelled{}.

\subsection{Gauge invisibility of a common datum}

\begin{lemma}[Gauge invisibility and the differential tax]
\label{lem:gauge}
A datum component \emph{common} to all providers---$\datum_a=\datum$ for every
$a$---is unobservable to any single-provider user and, being common, also cancels in
provider-versus-provider fusion, $\datum_{ab}=\datum_b-\datum_a=\bm 0$. It becomes a
real position error only against an \emph{absolute external tie}: a landmark surveyed
in a designated frame that does not carry the same gauge. The object a purely
multi-provider (fusion or hand-off) user pays is instead the \emph{differential} datum
$\datum_{ab}\neq\bm 0$ that two providers do not share---the interop tax, and the
quantity Section~\ref{sec:interephemeris} measures.
\end{lemma}

\begin{proof}
A single provider's internal range and PVT observables depend only on relative
geometry: distances, delays, and clock differences between points expressed in that
provider's own frame. These are invariant under the frame \emph{gauge}---the subspace
of similarity transforms that the provider's own observations cannot fix. For the
single-frame lunar datum this residual gauge is generically of dimension at most six
and collapses toward zero once lunar libration is observed and the scale becomes
observable \citep{baweja2026datum}; it is \emph{not} the full seven-parameter
similarity group. A common datum $\datum$ lying in that gauge changes no internal
observable, so it is invisible intra-provider. Now form the cross-provider difference
\eqref{eq:diffmodel}: a datum shared by both providers enters as
$\datum_{ab}=\datum_b-\datum_a=\bm 0$ and cancels identically---fusing or handing off
between two providers that share a gauge costs nothing. The common gauge becomes a
real position error only when a provider coordinate is tied to an \emph{absolute}
external reference that does not carry the same gauge, at which point it appears as
$\Jac(\bm{p})\datum$; adopting a designated realization (one provider's frame, or a
designated ILRF realization \citep{sosnica2025ilrf,altamimi2016itrf2014}) as the
shared datum removes it by construction. The tax that a purely multi-provider user
pays is therefore not the common gauge but the \emph{differential} datum
$\datum_{ab}\neq\bm 0$ of \eqref{eq:diffmodel}, the part the two providers do not
share.
\end{proof}

\noindent Lemma~\ref{lem:gauge} separates two objects that are easily conflated: a
\emph{common} gauge, which cancels in multi-provider fusion and bites only against an
absolute tie, and the \emph{differential} datum $\datum_{ab}$, which is invisible to
each provider in isolation but real the instant two providers are mixed. The interop
tax is the second. The lemma also identifies the mechanism by which a convention
helps---agreement that shrinks $\datum_{ab}$ toward a shared realization---and thereby
sets up the split.

\subsection{The reducible/irreducible split}

\begin{lemma}[Reducible/irreducible split]
\label{lem:split}
The cross-provider interop tax admits the measured affine split
$\reduc + \irr + \resid$, where $\reduc$ (the reducible part) is the body-common
Helmert rotation $\rot_{\mathrm{tie}}$, $\irr$ (the irreducible part) is the
body-specific Moon-excess rotation
$\rot_{\mathrm{exc}}=\rot_{\mathrm{moon}}-\rot_{\mathrm{tie}}$, and $\resid$ is the
non-Helmert residual. Adopting a common frame realization sets $\reduc\to0$ and
leaves $\irr$ and $\resid$ unchanged; only adopting a common designated dynamical
reference sets $\irr\to0$ (and drives $\resid\to0$ with it). The split is affine, not
orthogonal: $\rot_{\mathrm{tie}}$ and $\rot_{\mathrm{exc}}$ are in general
non-parallel.
\end{lemma}

\begin{proof}
Write the pairwise tax from \eqref{eq:diffmodel} as
$\Jac(\bm{p})\datum_{ab}+\resid_{ab}$. Split the recovered Moon rotation as
$\rot_{\mathrm{moon}}=\rot_{\mathrm{tie}}+\rot_{\mathrm{exc}}$, where
$\rot_{\mathrm{tie}}$ is the body-common frame-tie rotation estimated from the
planetary positions (Section~\ref{sec:interephemeris}) and
$\rot_{\mathrm{exc}}=\rot_{\mathrm{moon}}-\rot_{\mathrm{tie}}$ is the Moon-specific
remainder; set $\reduc=\Jac(\bm{p})\rot_{\mathrm{tie}}$,
$\irr=\Jac(\bm{p})\rot_{\mathrm{exc}}$, and keep $\resid_{ab}$ as a separate term.
This is a measured \emph{affine} decomposition, not an orthogonal one:
$\rot_{\mathrm{tie}}$ and $\rot_{\mathrm{exc}}$ need not be perpendicular---on the real
data they are $52^\circ$--$157^\circ$ apart (Section~\ref{sec:interephemeris})---so
the metre envelopes $\|\reduc\|$ and $\|\irr\|$ do not combine in quadrature to the
raw disagreement. What the proof establishes is the \emph{removability} structure. A
common frame realization is a single similarity transform applied to every
coordinate; it can cancel only a body-common rotation, so it sets $\reduc\to0$ and
leaves $\rot_{\mathrm{exc}}$ (hence $\irr$) and the residual $\resid_{ab}$
untouched---$\rot_{\mathrm{exc}}$ is by construction the part of the lunar rotation
absent from the planetary tie, and $\resid_{ab}$ is, in the aggregate least-squares
sense $\sum_i\Jac(\bm p_i)^{\!\top}\resid_{ab}(\bm p_i)=\bm 0$, the component of the
difference that no constant Helmert transform reaches. A common designated dynamical
reference replaces each provider's broadcast dynamics with one shared ephemeris, which
by construction makes the differenced dynamics vanish and sets $\irr\to0$ (and
$\resid_{ab}\to0$).
\end{proof}

\begin{remark}[The attribution is a modelling interpretation]
\label{rem:attribution}
Lemma~\ref{lem:split} establishes the removability structure---which convention
removes which part---and is \Proven{} at that level. The affine \emph{attribution} of
the recovered rotation into a frame-tie and a Moon-excess, and which physical cause we
assign to each, is a modelling choice. We
identify the body-common rotation with a removable ICRF \emph{frame-tie} (reducible)
and the Moon-specific excess with irreducible \emph{dynamics}. The convention-free,
attribution-independent content is the Moon-\emph{excess} rotation: it is by
construction the part of the lunar disagreement that does \emph{not} appear in the
planetary frame-tie, so it is unambiguously Moon-specific and therefore removable by
no whole-frame choice. The reader who declines the frame-tie interpretation may read
the split as ``planet-common rotation'' versus ``Moon-excess rotation'' and lose only
the word \emph{reducible}, not the irreducible floor. We restate this caveat wherever
the split is used (Sections~\ref{sec:interephemeris},~\ref{sec:budget},~%
\ref{sec:discussion}).
\end{remark}

\subsection{The consistency tolerance}

Let a mixed-provider user evaluate coordinates at lunocentric lever arm $\Rmoon$
(for a surface user, $\Rmoon=R_{\mathrm{Moon}}$; for the geocentric Moon state,
$\Rmoon$ is the Earth--Moon distance). The worst-case interop-tax position error from
a Helmert disagreement $\datum=(\ttrans,s,\rot)$ obeys a triangle bound.

\begin{theorem}[Consistency tolerance]
\label{thm:tolerance}
For a body-point at radius $\|\bm{p}\|\le\Rmoon$,
\begin{equation}
  \bigl|\Delta\xstate(\bm{p})\bigr|
  \;=\; \bigl|\Jac(\bm{p})\datum\bigr|
  \;\le\; |\ttrans| + |s|\,\Rmoon + |\rot|\,\Rmoon .
  \label{eq:triangle}
\end{equation}
Consequently, to hold the interop tax within a budget $B$, the cross-provider
consistency tolerances on the individual Helmert parameters are
\begin{equation}
  |\ttrans| \le \tau_t = B_{\mathrm{eff}},
  \qquad
  |s| \le \tau_s = \frac{B_{\mathrm{eff}}}{\Rmoon},
  \qquad
  |\rot| \le \tau_\theta = \frac{B_{\mathrm{eff}}}{\Rmoon},
  \label{eq:tau}
\end{equation}
with $B_{\mathrm{eff}}=B$ if the per-provider frame is realized exactly and
$B_{\mathrm{eff}}=\sqrt{\max(0,\,B^2-\sigma_{\mathrm{real}}^2)}$ otherwise, where
$\sigma_{\mathrm{real}}$ is the single-provider datum-realization standard deviation
from the companion frame-error covariance \citep{baweja2026datum}. At the lunar lever
arm the scale and rotation tolerances coincide at $B_{\mathrm{eff}}/\Rmoon$ and are
the binding (tightest) constraints; since $\Rmoon$ is large, an orientation error is
the cheapest way to exhaust the budget.
\end{theorem}

\begin{proof}
From \eqref{eq:pointjac},
$\Jac(\bm{p})\datum=\ttrans+s\bm{p}+\rot\times\bm{p}$. The triangle inequality gives
$|\Jac(\bm{p})\datum|\le|\ttrans|+|s|\,\|\bm{p}\|+|\rot\times\bm{p}|$, and
$|\rot\times\bm{p}|\le|\rot|\,\|\bm{p}\|\le|\rot|\,\Rmoon$ while
$|s|\,\|\bm{p}\|\le|s|\,\Rmoon$; substituting $\|\bm{p}\|\le\Rmoon$ yields
\eqref{eq:triangle}. For the tolerance: if a single Helmert parameter is nonzero, the
bound \eqref{eq:triangle} contributes exactly that term, and requiring it not to
exceed $B_{\mathrm{eff}}$ gives \eqref{eq:tau}; the translation term is dimensionless
in $\Rmoon$ while the scale and rotation terms carry the factor $\Rmoon$, so at
$\Rmoon\gg1$ the latter two bind first, at the common value $B_{\mathrm{eff}}/\Rmoon$.
The inflation $B_{\mathrm{eff}}=\sqrt{\max(0,B^2-\sigma_{\mathrm{real}}^2)}$ follows
from budgeting in quadrature: a frame the provider cannot even realize consumes
$\sigma_{\mathrm{real}}$ of the budget, leaving
$\sqrt{B^2-\sigma_{\mathrm{real}}^2}$ for cross-provider agreement.
\end{proof}

\noindent The triangle bound \eqref{eq:triangle} and the inversion form
\eqref{eq:tau} are \Proven{}. Each tolerance in \eqref{eq:tau} bounds a \emph{single}
Helmert parameter \emph{in isolation}: it is the largest that parameter may be if it
alone is nonzero. If translation, scale, and rotation saturate their tolerances
simultaneously, the triangle bound \eqref{eq:triangle} permits a combined error of up
to $3B_{\mathrm{eff}}$. A joint budget guarantee therefore requires allocating $B$
across the three terms---$B/3$ apiece for a hard split, or a quadrature allocation of
the expected magnitudes---rather than granting each parameter the full
$B_{\mathrm{eff}}$; \eqref{eq:tau} states the per-parameter ceilings, not a
simultaneous one. The $B_{\mathrm{eff}}$ inflation is \Modelled{}: it
depends on the companion CRLB $\sigma_{\mathrm{real}}$, itself a modelled magnitude,
and on treating realization and cross-provider errors as independent. When the
per-provider realization is taken as exact (as in the reported tolerance map,
Section~\ref{sec:budget}), $B_{\mathrm{eff}}=B$ and the tolerance is exact. The
quantity $\tau(B)$ is precisely what an interoperability specification must state to
convert a user PVT budget into a required per-provider frame-and-dynamics agreement:
it is the number a standard at the AD-5 level---still unreleased---must contain.

\subsection{The design-law corollary}

\begin{corollary}[Frame-versus-ephemeris design law]
\label{cor:designlaw}
Define the \emph{irreducible fraction} of the cross-provider rotation budget as the
linear ratio
\begin{equation}
  f_{\irr} \;:=\;
  \frac{\|\rot_{\mathrm{exc}}\|_{\mathrm{rms}}}
       {\|\rot_{\mathrm{tie}}\|_{\mathrm{rms}}+\|\rot_{\mathrm{exc}}\|_{\mathrm{rms}}} .
  \label{eq:irrfrac}
\end{equation}
Under the frame convention today's interoperability standards pursue---agree on frame
tags, format, and an ICRF realization, i.e.\ set $\reduc\to0$ but not $\irr\to0$---the
residual interop tax is carried by $\rot_{\mathrm{exc}}$ (plus the small non-Helmert
residual). If $f_{\irr}>\tfrac12$, the irreducible rotation exceeds the reducible one,
so a common frame/format convention leaves the \emph{dominant} share of the
cross-provider orientation disagreement in place; closing it requires additionally
mandating a common (or consistency-bounded) designated dynamical reference, which sets
$\irr\to0$.
\end{corollary}

\begin{proof}
Immediate from Lemma~\ref{lem:split}: adopting a common frame realization removes
$\reduc$ and leaves $\rot_{\mathrm{exc}}$, so the post-convention orientation tax is
$\|\rot_{\mathrm{exc}}\|$. The linear ratio \eqref{eq:irrfrac} exceeds $\tfrac12$
exactly when
$\|\rot_{\mathrm{exc}}\|_{\mathrm{rms}}>\|\rot_{\mathrm{tie}}\|_{\mathrm{rms}}$. Because
the split is affine rather than orthogonal (Lemma~\ref{lem:split}), we compare the
rotations directly and do not combine them in quadrature; the lever-arm-independent
rotations are the primary quantities, with metres following as worst-case envelopes at
a stated lever arm (Section~\ref{sec:budget}). Removing $\irr$ requires the differenced
dynamics to vanish, which a common designated ephemeris achieves by construction.
\end{proof}

\noindent Corollary~\ref{cor:designlaw} is a proven implication of the split; its
\emph{premise}---whether the irreducible fraction actually exceeds $\tfrac12$ for real
lunar providers---is an empirical question answered by the inter-ephemeris anchor of
the next section, where we measure an irreducible fraction of $0.687$ (the irreducible
rotation is about twice the reducible) for the representative provider set. That the
premise holds, and hence that a common format leaves a real floor, is the design law's
content; the specific magnitude of the floor is \Modelled{} (Section~\ref{sec:budget}).
